\section{Introdução teórica}
\label{sec:intro}

O experimento consiste em uma viga esbelta de aço com uma extremidade fixada a um
suporte rígido por porcas e parafusos e a outra extremidade livre, à qual podem estar
presas massas concentradas \cite{roteiroLab1}. A viga foi excitada de duas formas:
por um martelo modal aplicado a \SI{20}{\centi\meter} da extremidade fixa, com a
resposta medida por um acelerômetro a \SI{35.6}{\centi\meter}; e por um
\emph{shaker} acoplado à extremidade livre, em varredura senoidal, com força e
aceleração medidas por uma cabeça de impedância. O objetivo é extrair frequências
naturais e fatores de amortecimento das respostas medidas e compará-los com um modelo
analítico.

Esta seção desenvolve esse modelo. Como a fixação real não é perfeitamente rígida, e o
roteiro pede para avaliar qual condição de contorno descreve melhor a extremidade
fixa, são consideradas duas hipóteses: um engaste ideal e uma base elástica, com molas
torcional e transversal. As duas formulações são desenvolvidas em paralelo até a
equação característica e as formas modais não normalizadas
(seções~\ref{sec:eb}--\ref{sec:molas}). A normalização e a função de resposta em
frequência (FRF) seguem o mesmo procedimento nos dois casos, que só diferem nas
frequências naturais e nas formas modais (seções~\ref{sec:normalizacao} e
\ref{sec:frf}). Ao fim de cada etapa há um código em SymPy que refaz a dedução
simbólica e verifica o resultado.

%=====================================================================
\subsection{Modelo de Euler--Bernoulli e separação de variáveis}
\label{sec:eb}

Para uma viga esbelta, de seção e material uniformes, sob pequenos deslocamentos e
desprezando a inércia de rotação e a deformação por cisalhamento, o deslocamento
transversal $w(x,t)$ obedece à equação de Euler--Bernoulli
\cite{rao2019,meirovitch2001}
\begin{equation}
    EI\,\frac{\partial^4 w}{\partial x^4} + \rho A\,\frac{\partial^2 w}{\partial t^2}
    = f(x,t),
    \label{eq:eb}
\end{equation}
em que $E$ é o módulo de elasticidade, $I$ o momento de inércia de área da seção,
$\rho$ a massa específica, $A$ a área da seção e $f(x,t)$ a força distribuída por
unidade de comprimento.

Para obter as frequências naturais, toma-se $f = 0$ e procura-se uma solução da forma
$w(x,t) = W(x)\,g(t)$. Substituindo em \eqref{eq:eb} e separando as variáveis,
\begin{equation*}
    -\frac{\ddot g}{g} = \frac{EI}{\rho A}\,\frac{W''''}{W} = \omega^2,
\end{equation*}
pois o lado esquerdo depende só de $t$ e o direito só de $x$; logo, ambos são iguais
a uma constante, escrita como $\omega^2$. Resultam duas equações diferenciais
ordinárias:
\begin{subequations}
    \label{eq:separadas}
    \begin{align}
        \ddot g + \omega^2 g & = 0, \label{eq:temporal}                           \\
        W'''' - \beta^4 W    & = 0, \qquad \beta^4 = \frac{\rho A}{EI}\,\omega^2.
        \label{eq:espacial}
    \end{align}
\end{subequations}

A equação~\eqref{eq:temporal} é a de um sistema não amortecido de um grau de
liberdade, com solução harmônica $g = G\,e^{\iu\omega t}$. Para a parte espacial,
supondo $W = B\,e^{sx}$ em \eqref{eq:espacial},
\begin{equation*}
    (s^4 - \beta^4)\,B\,e^{sx} = 0 \quad\Longrightarrow\quad
    s_{1,2} = \pm\iu\beta, \qquad s_{3,4} = \pm\beta,
\end{equation*}
e a combinação das quatro exponenciais pode ser reescrita com funções trigonométricas
e hiperbólicas:
\begin{equation}
    W(x) = C_1\cos\beta x + C_2\sin\beta x + C_3\cosh\beta x + C_4\sinh\beta x.
    \label{eq:Wgeral}
\end{equation}

As constantes $C_1,\dots,C_4$ e os valores admissíveis de $\beta$ são determinados
pelas quatro condições de contorno, duas em cada extremidade. Com o parâmetro
adimensional $\lambda = \beta L$, as frequências naturais seguem de
\eqref{eq:espacial}:
\begin{equation}
    \omega_r = \beta_r^2\sqrt{\frac{EI}{\rho A}}
    = \frac{\lambda_r^2}{L^2}\sqrt{\frac{EI}{\rho A}},
    \qquad r = 1, 2, \dots
    \label{eq:omega}
\end{equation}

\lstinputlisting[caption={Separação de variáveis e solução geral da equação de
            Euler--Bernoulli.}, label={lst:solucao_geral}]{codigo/solucao_geral.py}

%=====================================================================
\subsection{Condições de contorno}
\label{sec:contorno}

A Figura~\ref{fig:modelos} mostra os dois modelos considerados. Em ambos, a
extremidade $x = L$ é livre e pode carregar uma massa concentrada $m$, como a massa
$m_{\mathrm{end}}$ da bancada descrita no roteiro \cite{roteiroLab1}. Os ensaios
analisados neste relatório foram feitos sem essa massa, isto é, com $m = 0$. A
formulação geral é mantida porque permite estimar, na Seção~\ref{sec:massa}, quanto as
frequências naturais mudariam com a massa, e porque uma massa pequena também
representa a cabeça de impedância presa à ponta no ensaio com \emph{shaker}. Os
modelos diferem apenas na extremidade $x = 0$.

\begin{figure}[htbp]
    \centering
    \begin{subfigure}{\textwidth}
        \centering
        \begin{tikzpicture}[>=Latex, scale=1.1]
            % engaste
            \fill[pattern=north east lines] (-0.35,-0.9) rectangle (0,0.9);
            \draw[thick] (0,-0.9) -- (0,0.9);
            % viga
            \draw[fill=gray!20] (0,-0.1) rectangle (6,0.1);
            \draw[dash dot, gray] (0,0) -- (6,0);
            % massa na ponta
            \draw[fill=gray!60] (6,0) circle (0.28);
            \node at (6,0.55) {$m$};
            % eixos
            \draw[->] (6.4,0) -- (7.2,0) node[right] {$x$};
            \draw[->] (3.6,0.1) -- (3.6,0.9) node[above] {$w(x,t)$};
            % força
            \draw[->, thick] (2.0,1.0) -- (2.0,0.12);
            \node[above] at (2.0,1.0) {$f(x,t)$};
            % propriedades e cota
            \node at (4.6,-0.4) {$EI,\ \rho A$};
            \draw[|<->|] (0,-1.15) -- (6,-1.15) node[midway, fill=white] {$L$};
            \node[above] at (-0.17,0.9) {$x=0$};
        \end{tikzpicture}
        \caption{Base engastada: $w(0,t) = 0$ e $w'(0,t) = 0$.}
        \label{fig:modelo_engaste}
    \end{subfigure}

    \vspace{0.6cm}

    \begin{subfigure}{\textwidth}
        \centering
        \begin{tikzpicture}[>=Latex, scale=1.1]
            % parede da mola torcional
            \fill[pattern=north east lines] (-1.35,-0.6) rectangle (-1.05,0.6);
            \draw[thick] (-1.05,-0.6) -- (-1.05,0.6);
            % viga
            \draw[fill=gray!20] (0,-0.1) rectangle (6,0.1);
            \draw[dash dot, gray] (0,0) -- (6,0);
            % mola torcional (espiral centrada em x = 0, termina em 180 graus)
            \fill[white] (0,0) circle (0.4);
            \draw[thick, domain=0:3*pi, samples=150, variable=\t]
            plot ({\t r}:{0.04*\t});
            \draw[thick] (-0.377,0) -- (-1.05,0);
            \fill (0,0) circle (0.045);
            \node at (-0.55,0.45) {$k_R$};
            % mola transversal
            \draw[thick] (0.6,-0.1) -- (0.6,-0.35);
            \draw[thick, decorate,
                decoration={zigzag, segment length=5pt, amplitude=4pt}]
            (0.6,-0.35) -- (0.6,-1.05);
            \draw[thick] (0.6,-1.05) -- (0.6,-1.2);
            \fill[pattern=north east lines] (0.25,-1.45) rectangle (0.95,-1.2);
            \draw[thick] (0.25,-1.2) -- (0.95,-1.2);
            \node[right] at (0.75,-0.7) {$k_L$};
            % massa na ponta
            \draw[fill=gray!60] (6,0) circle (0.28);
            \node at (6,0.55) {$m$};
            % eixos
            \draw[->] (6.4,0) -- (7.2,0) node[right] {$x$};
            \draw[->] (3.6,0.1) -- (3.6,0.9) node[above] {$w(x,t)$};
            % força
            \draw[->, thick] (2.0,1.0) -- (2.0,0.12);
            \node[above] at (2.0,1.0) {$f(x,t)$};
            % propriedades e cota
            \node at (4.6,-0.4) {$EI,\ \rho A$};
            \draw[|<->|] (0,-1.75) -- (6,-1.75) node[midway, fill=white] {$L$};
        \end{tikzpicture}
        \caption{Base elástica: mola torcional $k_R$ e mola transversal $k_L$ em
            $x = 0$.}
        \label{fig:modelo_molas}
    \end{subfigure}
    \caption{Modelos da viga com massa concentrada na extremidade livre.}
    \label{fig:modelos}
\end{figure}

Os esforços internos seguem a convenção da Figura~\ref{fig:convencao}a: momento
fletor $M = EI\,w''$ e esforço cortante $Q = -EI\,w'''$, positivos como indicado na
face de normal $+x$ e opostos na face de normal $-x$. Com essa convenção, o equilíbrio
do elemento dá $M' = -Q$.

\begin{figure}[htbp]
    \centering
    \begin{tikzpicture}[>=Latex, scale=1.0]
        % ---------- (a) elemento
        \begin{scope}
            \draw[fill=gray!20] (0,-0.3) rectangle (1.6,0.3);
            \draw[->, thick] (-0.15,0.65) -- (-0.15,-0.65) node[below] {$Q$};
            \draw[->, thick] (1.75,-0.65) -- (1.75,0.65) node[above] {$Q + \mathrm{d}Q$};
            \draw[->, thick] (0,0) ++(215:0.6) arc (215:145:0.6);
            \node[left] at (-0.6,0) {$M$};
            \draw[->, thick] (1.6,0) ++(-35:0.6) arc (-35:35:0.6);
            \node[right] at (2.2,0) {$M + \mathrm{d}M$};
            \node at (0.8,-1.8) {(a)};
        \end{scope}
        % ---------- (b) base com molas
        \begin{scope}[xshift=6.6cm]
            \draw[fill=gray!20] (0,-0.1) rectangle (1.8,0.1);
            \draw[->, thick] (0,0.75) -- (0,0.12);
            \node[above] at (0,0.75) {$k_L\,w(0)$};
            \draw[->, thick] (0,0) ++(215:0.5) arc (215:145:0.5);
            \node[left] at (-0.5,0) {$k_R\,w'(0)$};
            \node at (0.9,-1.8) {(b)};
        \end{scope}
        % ---------- (c) massa na ponta
        \begin{scope}[xshift=11.6cm]
            \draw[fill=gray!60] (0,0) circle (0.3);
            \node at (0,0) {$m$};
            \draw[->, thick] (0,-0.32) -- (0,-1.0) node[right] {$Q(L)$};
            \draw[->, thick, dashed] (0.55,0) -- (0.55,0.8) node[above] {$m\,\ddot w(L)$};
            \node at (0,-1.8) {(c)};
        \end{scope}
    \end{tikzpicture}
    \caption{(a) Convenção de sinais dos esforços internos em um elemento $\mathrm{d}x$.
        (b) Ações das molas sobre a viga em $x=0$ para $w(0)>0$ e $w'(0)>0$.
        (c) Diagrama de corpo livre da massa na ponta.}
    \label{fig:convencao}
\end{figure}

\subsubsection{Extremidade livre com massa concentrada ($x = L$)}

Tratando a massa como uma partícula (sem inércia de rotação), ela não transmite
momento, e a única força que recebe é a cortante da viga
(Fig.~\ref{fig:convencao}c). Assim,
\begin{subequations}
    \label{eq:cc_ponta}
    \begin{align}
        M(L,t) = EI\,w''(L,t) = 0
         & \quad\Longrightarrow\quad W''(L) = 0, \label{eq:cc_ponta_M}           \\
        Q(L,t) = -EI\,w'''(L,t) = -m\,\ddot w(L,t)
         & \quad\Longrightarrow\quad W'''(L) = -\frac{m}{\rho A}\,\beta^4\,W(L),
        \label{eq:cc_ponta_Q}
    \end{align}
\end{subequations}
onde se usou $\ddot g = -\omega^2 g$, que dá $EI\,W'''(L) = -m\,\omega^2 W(L)$, e
\eqref{eq:espacial}. Definindo a razão entre a
massa da ponta e a massa da viga,
\begin{equation}
    \mu = \frac{m}{\rho A L},
    \label{eq:mu}
\end{equation}
e dividindo \eqref{eq:cc_ponta_Q} por $\beta^3$, a condição de cortante fica
$W'''(L)/\beta^3 = -\mu\lambda\,W(L)$. Essas duas condições valem para os dois
modelos.

\subsubsection{Estrutura comum do problema de autovalor}
\label{sec:estrutura}

Em cada modelo, as duas condições em $x = 0$ eliminam $C_3$ e $C_4$ de
\eqref{eq:Wgeral}, e a forma modal passa a depender só de $C_1$ e $C_2$:
\begin{equation}
    W(x) = C_1\,\phi_1(\beta x) + C_2\,\phi_2(\beta x),
    \label{eq:W_phi}
\end{equation}
em que $\phi_1$ e $\phi_2$ são funções conhecidas para cada modelo. Substituindo
\eqref{eq:W_phi} nas condições \eqref{eq:cc_ponta} e indicando por $\phi_j'$ a derivada
em relação ao argumento $\zeta = \beta x$, obtém-se o sistema homogêneo
\begin{equation}
    \underbrace{\begin{bmatrix}
            \phi_1''(\lambda)                                & \phi_2''(\lambda) \\[2pt]
            \phi_1'''(\lambda) + \mu\lambda\,\phi_1(\lambda) &
            \phi_2'''(\lambda) + \mu\lambda\,\phi_2(\lambda)
        \end{bmatrix}}_{\mathbf{B}(\lambda)}
    \begin{Bmatrix} C_1 \\ C_2 \end{Bmatrix}
    = \begin{Bmatrix} 0 \\ 0 \end{Bmatrix}.
    \label{eq:sistema}
\end{equation}
A solução não trivial exige $\det\mathbf{B}(\lambda) = 0$, que é a
equação característica. Suas raízes $\lambda_r$ dão as frequências naturais por
\eqref{eq:omega}. Para cada raiz, a primeira linha de \eqref{eq:sistema} fixa a razão
entre as constantes,
\begin{equation}
    \gamma_r = \frac{C_2}{C_1} = -\frac{\phi_1''(\lambda_r)}{\phi_2''(\lambda_r)},
    \label{eq:gamma}
\end{equation}
e a forma modal fica determinada a menos de uma constante:
\begin{equation}
    W_r(x) = C_r\,F_r(x), \qquad
    F_r(x) = \phi_1(\beta_r x) + \gamma_r\,\phi_2(\beta_r x).
    \label{eq:Fr}
\end{equation}

\lstinputlisting[caption={Condições de contorno na extremidade com massa
            concentrada, em forma adimensional.}, label={lst:contorno}]{codigo/contorno_ponta.py}

%=====================================================================
\subsection{Caso 1: base engastada}
\label{sec:engaste}

No engaste ideal (Fig.~\ref{fig:modelo_engaste}), o deslocamento e a rotação são nulos:
\begin{subequations}
    \label{eq:cc_engaste}
    \begin{align}
        w(0,t) = 0  & \quad\Longrightarrow\quad W(0) = C_1 + C_3 = 0,           \\
        w'(0,t) = 0 & \quad\Longrightarrow\quad W'(0) = \beta\,(C_2 + C_4) = 0,
    \end{align}
\end{subequations}
de modo que $C_3 = -C_1$ e $C_4 = -C_2$. A forma \eqref{eq:W_phi} fica
\begin{equation}
    W(x) = C_1\left(\cos\beta x - \cosh\beta x\right)
    + C_2\left(\sin\beta x - \sinh\beta x\right),
    \label{eq:W_engaste}
\end{equation}
ou seja, $\phi_1(\zeta) = \cos\zeta - \cosh\zeta$ e
$\phi_2(\zeta) = \sin\zeta - \sinh\zeta$. As derivadas necessárias são
\begin{align*}
    \phi_1'' & = -(\cos\zeta + \cosh\zeta), & \phi_1''' & = \sin\zeta - \sinh\zeta,    \\
    \phi_2'' & = -(\sin\zeta + \sinh\zeta), & \phi_2''' & = -(\cos\zeta + \cosh\zeta),
\end{align*}
e os elementos da matriz de \eqref{eq:sistema} são
\begin{equation}
    \begin{aligned}
        A_{11} & = -(\cos\lambda + \cosh\lambda),                                          &
        A_{21} & = \sin\lambda - \sinh\lambda + \mu\lambda(\cos\lambda - \cosh\lambda),      \\
        A_{12} & = -(\sin\lambda + \sinh\lambda),                                          &
        A_{22} & = -(\cos\lambda + \cosh\lambda) + \mu\lambda(\sin\lambda - \sinh\lambda).
    \end{aligned}
    \label{eq:A_engaste}
\end{equation}

Expandindo o determinante e usando $\cos^2\lambda + \sin^2\lambda = 1$ e
$\cosh^2\lambda - \sinh^2\lambda = 1$, chega-se a
$\det\mathbf{B} = 2\left[1 + \cos\lambda\cosh\lambda
        + \mu\lambda(\cos\lambda\sinh\lambda - \sin\lambda\cosh\lambda)\right]$, e a equação
característica é
\begin{equation}
    1 + \cos\lambda\cosh\lambda
    + \mu\lambda\left(\cos\lambda\sinh\lambda - \sin\lambda\cosh\lambda\right) = 0.
    \label{eq:car_engaste}
\end{equation}
A equação é transcendental e suas raízes $\lambda_r$ precisam ser obtidas
numericamente. Sem massa na ponta ($\mu = 0$), ela se reduz à equação clássica da viga
engastada-livre, $1 + \cos\lambda\cosh\lambda = 0$, com
$\lambda_r = 1{,}8751;\ 4{,}6941;\ 7{,}8548;\ \dots$ \cite{rao2019}. A massa na ponta
reduz esses valores.

Por \eqref{eq:gamma}, a forma modal não normalizada do modo $r$ é
\begin{equation}
    F_r(x) = \cos\beta_r x - \cosh\beta_r x
    - \frac{\cos\lambda_r + \cosh\lambda_r}{\sin\lambda_r + \sinh\lambda_r}
    \left(\sin\beta_r x - \sinh\beta_r x\right).
    \label{eq:F_engaste}
\end{equation}

\lstinputlisting[caption={Caso engastado: matriz $\mathbf{A}$, equação característica e
            forma modal.}, label={lst:engaste}]{codigo/caso_engastado.py}

%=====================================================================
\subsection{Caso 2: base elástica}
\label{sec:molas}

No segundo modelo (Fig.~\ref{fig:modelo_molas}), a fixação em $x=0$ é representada
por uma mola torcional de rigidez $k_R$, que resiste à rotação, e uma mola transversal
de rigidez $k_L$, que resiste ao deslocamento. Para $w(0) > 0$ e $w'(0) > 0$, as
molas aplicam sobre a viga uma força $k_L w(0)$ para baixo e um momento
$k_R w'(0)$ horário (Fig.~\ref{fig:convencao}b). Pela convenção da
Figura~\ref{fig:convencao}a, essas ações correspondem a
\begin{subequations}
    \label{eq:cc_molas}
    \begin{align}
        M(0,t) & = EI\,w''(0,t) = k_R\,w'(0,t), \label{eq:cc_molas_M}  \\
        Q(0,t) & = -EI\,w'''(0,t) = k_L\,w(0,t). \label{eq:cc_molas_Q}
    \end{align}
\end{subequations}
Essas condições também são as condições naturais do princípio de Hamilton quando as
molas armazenam a energia $\tfrac12 k_R w'(0)^2 + \tfrac12 k_L w(0)^2$
\cite{meirovitch2001}. Substituindo \eqref{eq:Wgeral},
\begin{align*}
    EI\beta^2\left(-C_1 + C_3\right)  & = k_R\,\beta\left(C_2 + C_4\right), \\
    -EI\beta^3\left(-C_2 + C_4\right) & = k_L\left(C_1 + C_3\right).
\end{align*}
Definindo as rigidezes adimensionais
\begin{equation}
    \alpha_R = \frac{k_R}{EI\beta} = \frac{\kappa_R}{\lambda}, \qquad
    \alpha_L = \frac{k_L}{EI\beta^3} = \frac{\kappa_L}{\lambda^3}, \qquad
    \kappa_R = \frac{k_R L}{EI}, \qquad \kappa_L = \frac{k_L L^3}{EI},
    \label{eq:alphas}
\end{equation}
as condições ficam $C_3 = C_1 + \alpha_R(C_2 + C_4)$ e
$C_4 = C_2 - \alpha_L(C_1 + C_3)$. Resolvendo esse par de equações para $C_3$ e $C_4$:
\begin{equation}
    C_3 = \frac{(1 - \alpha_L\alpha_R)\,C_1 + 2\alpha_R\,C_2}{1 + \alpha_L\alpha_R},
    \qquad
    C_4 = \frac{-2\alpha_L\,C_1 + (1 - \alpha_L\alpha_R)\,C_2}{1 + \alpha_L\alpha_R}.
    \label{eq:C3C4_molas}
\end{equation}
Os parâmetros $\kappa_R$ e $\kappa_L$ descrevem a fixação e não dependem da
frequência; já $\alpha_R$ e $\alpha_L$ dependem de $\lambda$, porque a rigidez relativa
da viga muda com o comprimento de onda do modo. Como $1 + \alpha_L\alpha_R > 0$, esse
fator pode ser absorvido nas constantes. Redefinindo $C_1$ e $C_2$, a forma
\eqref{eq:W_phi} vale com
\begin{subequations}
    \label{eq:phi_molas}
    \begin{align}
        \phi_1(\zeta) & = (1 + \alpha_L\alpha_R)\cos\zeta
        + (1 - \alpha_L\alpha_R)\cosh\zeta - 2\alpha_L\sinh\zeta, \\
        \phi_2(\zeta) & = (1 + \alpha_L\alpha_R)\sin\zeta
        + (1 - \alpha_L\alpha_R)\sinh\zeta + 2\alpha_R\cosh\zeta.
    \end{align}
\end{subequations}
Os elementos da matriz $\mathbf{B}$ de \eqref{eq:sistema} são, então,
\begin{align*}
    B_{11} & = -(1 + \alpha_L\alpha_R)\cos\lambda + (1 - \alpha_L\alpha_R)\cosh\lambda
    - 2\alpha_L\sinh\lambda,                                                           \\
    B_{12} & = -(1 + \alpha_L\alpha_R)\sin\lambda + (1 - \alpha_L\alpha_R)\sinh\lambda
    + 2\alpha_R\cosh\lambda,                                                           \\
    B_{21} & = (1 + \alpha_L\alpha_R)\sin\lambda + (1 - \alpha_L\alpha_R)\sinh\lambda
    - 2\alpha_L\cosh\lambda + \mu\lambda\,\phi_1(\lambda),                             \\
    B_{22} & = -(1 + \alpha_L\alpha_R)\cos\lambda + (1 - \alpha_L\alpha_R)\cosh\lambda
    + 2\alpha_R\sinh\lambda + \mu\lambda\,\phi_2(\lambda).
\end{align*}
O determinante tem a forma $\det\mathbf{B} = 2(1 + \alpha_L\alpha_R)\,\Delta(\lambda)$,
e, como o primeiro fator não se anula, a equação característica é
$\Delta(\lambda) = 0$:
\begin{equation}
    \begin{split}
        \Delta(\lambda) ={} & 1 + \alpha_L\alpha_R
        + (\alpha_L\alpha_R - 1)\cos\lambda\cosh\lambda                      \\
                            & + (\alpha_L - \alpha_R)\cos\lambda\sinh\lambda
        - (\alpha_L + \alpha_R)\sin\lambda\cosh\lambda                       \\
                            & + \mu\lambda\Big[(1 - \alpha_L\alpha_R)
        (\sin\lambda\cosh\lambda - \cos\lambda\sinh\lambda)                  \\
                            & \qquad\quad - 2\alpha_L\sin\lambda\sinh\lambda
            - 2\alpha_R\cos\lambda\cosh\lambda\Big] = 0.
    \end{split}
    \label{eq:car_molas}
\end{equation}
Substituindo \eqref{eq:alphas} e multiplicando por $\lambda^4$, obtém-se uma função
apenas de $\lambda$ para valores dados de $\kappa_R$, $\kappa_L$ e $\mu$, cujas raízes
são calculadas numericamente. Há dois limites que servem de verificação:
\begin{itemize}
    \item para $\kappa_R, \kappa_L \to \infty$ (molas infinitamente rígidas), dividir
          \eqref{eq:car_molas} por $\alpha_L\alpha_R$ e fazer $\alpha_L,\alpha_R\to\infty$
          recupera exatamente a equação do engaste, \eqref{eq:car_engaste};
    \item para $\kappa_R = \kappa_L = \mu = 0$, resta $1 - \cos\lambda\cosh\lambda = 0$,
          a equação da viga livre-livre.
\end{itemize}
Entre esses limites, as frequências naturais crescem com as rigidezes da base. Isso
permite ajustar $\kappa_R$ e $\kappa_L$ às frequências medidas e avaliar quanto a
fixação real se afasta do engaste ideal.

A forma modal não normalizada segue de \eqref{eq:gamma} e \eqref{eq:Fr}:
\begin{equation}
    F_r(x) = \phi_1(\beta_r x) - \frac{B_{11}(\lambda_r)}{B_{12}(\lambda_r)}\,
    \phi_2(\beta_r x),
    \label{eq:F_molas}
\end{equation}
com $\alpha_R$ e $\alpha_L$ avaliados em $\lambda = \lambda_r$.

\lstinputlisting[caption={Caso com base elástica: constantes $C_3$ e $C_4$, matriz
            $\mathbf{B}$, equação característica e limite engastado.},
    label={lst:molas}]{codigo/caso_molas.py}

%=====================================================================
\subsection{Normalização dos modos}
\label{sec:normalizacao}

As formas modais \eqref{eq:F_engaste} e \eqref{eq:F_molas} estão definidas a menos de
uma constante: a amplitude física da resposta é o produto dessa constante pela
amplitude de $g(t)$. Para fixar $C_r$, usa-se a ortogonalidade dos modos.
Multiplicando \eqref{eq:espacial} do modo $r$ por $W_s$, integrando duas vezes por
partes em $[0, L]$ e usando as condições de contorno \eqref{eq:cc_ponta} e
\eqref{eq:cc_molas}, obtém-se
\begin{multline}
    \int_0^L EI\,W_r''W_s''\,\mathrm{d}x + k_R\,W_r'(0)W_s'(0) + k_L\,W_r(0)W_s(0) \\
    = \omega_r^2\left[\int_0^L \rho A\,W_rW_s\,\mathrm{d}x + m\,W_r(L)W_s(L)\right].
    \label{eq:ortog_deriv}
\end{multline}
No caso engastado, os termos das molas desaparecem porque $W(0) = W'(0) = 0$. O lado
esquerdo é simétrico em $r$ e $s$. Trocando os índices e subtraindo, resulta
$(\omega_r^2 - \omega_s^2) = 0$, de modo que, para frequências distintas, os
modos são ortogonais em relação à massa, que inclui a da ponta. Escolhe-se então
$C_r$ de modo que
\begin{equation}
    \int_0^L \rho A\,W_r W_s\,\mathrm{d}x + m\,W_r(L)\,W_s(L) = \delta_{rs},
    \label{eq:ortonormal}
\end{equation}
e, por \eqref{eq:ortog_deriv}, o lado esquerdo daquela equação vale $\omega_r^2\,\delta_{rs}$.
Com $W_r = C_r F_r$ e seção e densidade constantes,
\begin{equation}
    C_r^2\left(\rho A\int_0^L F_r^2\,\mathrm{d}x + m\,F_r^2(L)\right) = 1,
\end{equation}
e os modos normalizados são
\begin{equation}
    W_r(x) = \frac{F_r(x)}{\left(\rho A\displaystyle\int_0^L F_r^2\,\mathrm{d}x
        + m\,F_r^2(L)\right)^{1/2}}.
    \label{eq:Wr_norm}
\end{equation}
A expressão é a mesma para os dois modelos: basta usar $F_r$ de
\eqref{eq:F_engaste} ou de \eqref{eq:F_molas} e as raízes $\lambda_r$ correspondentes.

\lstinputlisting[caption={Normalização pela massa e verificação da ortogonalidade de
            massa e de rigidez para os dois modelos.}, label={lst:normalizacao}]{codigo/normalizacao.py}

%=====================================================================
\subsection{Resposta forçada e função de resposta em frequência}
\label{sec:frf}

A resposta forçada é escrita como superposição modal,
\begin{equation}
    w(x,t) = \sum_{r=1}^{\infty} W_r(x)\,\eta_r(t),
    \label{eq:expansao}
\end{equation}
em que $\eta_r$ são as coordenadas modais. Substituindo \eqref{eq:expansao} em
\eqref{eq:eb}, multiplicando por $W_s$, integrando ao longo da viga e usando
\eqref{eq:ortonormal} e \eqref{eq:ortog_deriv}, as equações se desacoplam. Acrescenta-se
um amortecimento viscoso modal $\zeta_r$ (hipótese de amortecimento leve e
proporcional, \cite{ewins2000}):
\begin{equation}
    \ddot\eta_r + 2\zeta_r\omega_r\,\dot\eta_r + \omega_r^2\,\eta_r = N_r(t),
    \qquad N_r(t) = \int_0^L W_r(x)\,f(x,t)\,\mathrm{d}x.
    \label{eq:modal}
\end{equation}
Sem forçamento e sem amortecimento, \eqref{eq:modal} reduz-se a
$\ddot\eta_r + \omega_r^2\eta_r = 0$, a mesma equação \eqref{eq:temporal}. Cada
coordenada modal se comporta, portanto, como um sistema de um grau de liberdade.
A solução é a soma de uma parcela homogênea, que depende das condições iniciais e
decai com o amortecimento, e de uma parcela particular, que é a resposta em regime
permanente.

Para uma força harmônica pontual aplicada em $x = x_i$,
$f(x,t) = F_0\,\delta(x - x_i)\,e^{\iu\Omega t}$, a força generalizada é
\begin{equation}
    N_r(t) = F_0\,e^{\iu\Omega t}\int_0^L \delta(x - x_i)\,W_r(x)\,\mathrm{d}x
    = F_0\,W_r(x_i)\,e^{\iu\Omega t}.
    \label{eq:Nr}
\end{equation}

A solução particular de \eqref{eq:modal} é, então
\begin{equation}
    \eta_{r,p}(t) = \frac{F_0\,W_r(x_i)}{\omega_r^2}\,G_r(\iu\Omega)\,e^{\iu\Omega t},
    \qquad
    G_r(\iu\Omega) = \frac{1}{1 - \left(\Omega/\omega_r\right)^2
        + 2\iu\zeta_r\,\Omega/\omega_r},
    \label{eq:eta_p}
\end{equation}
com módulo e fase
\begin{equation}
    |G_r| = \frac{1}{\sqrt{\left[1 - (\Omega/\omega_r)^2\right]^2
            + \left(2\zeta_r\Omega/\omega_r\right)^2}},
    \qquad
    \angle G_r = -\arctan \left[ \frac{2\zeta_r\,\Omega/\omega_r}{1 - (\Omega/\omega_r)^2} \right].
\end{equation}
Desprezando a parcela homogênea, que já decaiu no regime permanente, a resposta é
\begin{equation}
    w(x,t) = F_0\,e^{\iu\Omega t}\sum_{r=1}^{\infty}
    \frac{W_r(x)\,W_r(x_i)}{\omega_r^2 - \Omega^2 + 2\iu\zeta_r\,\omega_r\Omega}.
    \label{eq:resposta}
\end{equation}

A FRF entre um ponto de saída $x_e$ e um ponto de entrada $x_i$ é a razão entre a
resposta em $x_e$ e a força em $x_i$. Para deslocamento (receptância), de
\eqref{eq:resposta},
\begin{equation}
    H(x_e, x_i, \Omega) = \frac{w(x_e,t)}{f(x_i,t)}
    = \sum_{r=1}^{\infty}
    \frac{W_r(x_e)\,W_r(x_i)}{\omega_r^2 - \Omega^2 + 2\iu\zeta_r\,\omega_r\Omega}.
    \label{eq:H}
\end{equation}
A expressão é simétrica em $x_e$ e $x_i$ (reciprocidade). Para o cálculo, a série é
truncada em um número de modos que cubra com folga a faixa de frequências medida.
Como os ensaios medem aceleração, a comparação com os dados experimentais usa a
acelerância, obtida de \eqref{eq:H}:
\begin{equation}
    \mathcal{A}(x_e, x_i, \Omega) = \frac{\ddot w(x_e,t)}{f(x_i,t)}
    = -\Omega^2\,H(x_e, x_i, \Omega).
    \label{eq:acelerancia}
\end{equation}
No ensaio com martelo, $x_i = \SI{0.20}{\meter}$ e $x_e = \SI{0.356}{\meter}$. No
ensaio com \emph{shaker}, a cabeça de impedância mede força e aceleração no mesmo
ponto, então a FRF é pontual, com $x_i = x_e = L$. Em ambos os casos, a FRF analítica é
obtida com as frequências $\omega_r$ e os modos $W_r$ de um dos dois modelos de base e
com os fatores $\zeta_r$ identificados experimentalmente.

\lstinputlisting[caption={Equação modal, resposta harmônica e FRF analítica
            (receptância e acelerância).}, label={lst:frf}]{codigo/frf.py}
